%!TEX root = ../../note.tex
\subsection{The Natural Numbers: Induction and Well-Ordering}

Two properties of $\N$ are used throughout the course. Induction lets us
prove a statement for every $n$ by passing from $n$ to $n+1$; well-ordering
lets us pick the least element of a set of natural numbers. We take the first
as a principle and derive the second from it.

\subsubsection{Mathematical induction}
\begin{axiom}[Principle of Mathematical Induction]
    Let $P(n)$ be a statement depending on $n\in\N$. If
    \begin{enumerate}[label=(\arabic*)]
        \item $P(1)$ is true, and
        \item for every $k\in\N$, $P(k)$ implies $P(k+1)$,
    \end{enumerate}
    then $P(n)$ is true for every $n\in\N$.
\end{axiom}

Thus an induction proof has two parts: the \term{base case} $P(1)$, and the
\term{inductive step}, in which we fix an arbitrary $k$, assume $P(k)$ (the
\emph{inductive hypothesis}), and deduce $P(k+1)$.

\begin{example}
    $1+2+\cdots+n=\dfrac{n(n+1)}{2}$ for every $n\in\N$.
\end{example}

\begin{proof}
    For $n=1$ both sides equal $1$. If the formula holds for $k$, then
    \begin{equation*}
        1+\cdots+k+(k+1)
        = \frac{k(k+1)}{2}+(k+1)
        = \frac{k(k+1)+2(k+1)}{2}
        = \frac{(k+1)(k+2)}{2},
    \end{equation*}
    which is the formula for $k+1$.
\end{proof}

The same principle, stated for sets, reads as follows.

\begin{proposition}\label{prop:induction-sets}
    If $S\subseteq\N$, $1\in S$, and $k\in S\Rightarrow k+1\in S$, then
    $S=\N$.
\end{proposition}

\begin{proof}
    Apply induction to the statement $P(n)$: ``$n\in S$''.
\end{proof}

\subsubsection{The well-ordering principle}
\begin{theorem}[Well-Ordering Principle]\label{thm:wop}
    Every nonempty subset $S$ of $\N$ has a least element.
\end{theorem}

Suppose $S$ had no least element. We show by induction that then $S$
contains none of $1,2,\ldots,n$, for every $n$; hence $S=\emptyset$.

\begin{proof}
    Assume $S\neq\emptyset$ has no least element, and let
    \begin{equation*}
        T=\set{n\in\N : \set{1,2,\ldots,n}\cap S=\emptyset}.
    \end{equation*}
    \emph{$1\in T$}: otherwise $1\in S$, and $1$, being the least natural
    number, would be the least element of $S$.

    \emph{$k\in T\Rightarrow k+1\in T$}: if $\set{1,\ldots,k}\cap S=\emptyset$
    but $\set{1,\ldots,k+1}\cap S\neq\emptyset$, then $k+1\in S$ while no
    smaller number lies in $S$, so $k+1$ would be the least element of $S$.

    By Proposition~\ref{prop:induction-sets}, $T=\N$. Hence every $n\in\N$
    satisfies $n\notin S$, i.e.\ $S=\emptyset$, a contradiction.
\end{proof}

Well-ordering gives the method of the \term{minimal counterexample}: if
$P(n)$ failed for some $n$, the set $S=\set{n\in\N:\neg P(n)}$ would be
nonempty and would have a least element $k$. Then $P(k-1)$ holds by
minimality, and one shows that $P(k-1)$ forces $P(k)$ --- a contradiction.

\begin{example}
    Using well-ordering, prove that for every $n\in\N$
    \begin{equation}\label{eq:sum-squares}
        1^2 + 2^2 + \cdots + n^2 = \frac{n(n + 1)(2n + 1)}{6}.
    \end{equation}
\end{example}

\begin{proof}
    Let $P(n)$ be the statement \eqref{eq:sum-squares}, and suppose that
    the set of counterexamples
    \begin{equation*}
        S=\set{n \in \N : \neg P(n)}
    \end{equation*}
    is nonempty. By
    Theorem~\ref{thm:wop}, $S$ has a least element $k$. Since $P(1)$ holds,
    $k>1$, so $k-1\in\N$, and $P(k-1)$ holds because $k-1<k$:
    \begin{equation*}
        1^2 + \cdots + (k - 1)^2 = \frac{(k - 1)k(2k - 1)}{6}.
    \end{equation*}
    Adding $k^2$,
    \begin{align*}
        1^2 + \cdots + k^2
        &= \frac{k\bigl((k-1)(2k-1)+6k\bigr)}{6} \\
        &= \frac{k(2k^2+3k+1)}{6}
         = \frac{k(k+1)(2k+1)}{6},
    \end{align*}
    so $P(k)$ holds, contradicting $k\in S$. Hence $S=\emptyset$.
\end{proof}

\begin{remark}
    The argument needs the \emph{least element} $k$ of $S$, not just a
    greatest lower bound: we must know that $k\in S$, i.e.\ that $P(k)$
    fails. For a nonempty subset of $\N$ this is exactly what
    well-ordering provides; for a subset of $\R$ the infimum need not belong
    to the set (\S\ref{sec:bounds}).
\end{remark}
